\chapter{Scientific scope and limitations}

\eyebrow{what the numbers can and cannot claim}

\vspace{4pt}
This chapter states the boundary between what the measurements of Apex establish
and what remains open.

\section{What is validated}

\begin{itemize}[leftmargin=1.3em,itemsep=3pt]
  \item \textbf{Fat segmentation} is validated (IoU $\sim$0.92) and reliable, once
        the background has been removed so the model sees the conditions it was
        trained on.
  \item \textbf{Verifiable pairing} is guaranteed by construction: every image
        carries its \tele{carcass\_id} from the moment of capture or import.
  \item \textbf{Conformation convexity} is an objective, reproducible measurement
        of the profile, invariant to pose and lighting; its correspondence to
        anatomy is verifiable by eye.
\end{itemize}

\section{What is not validated}

The two \emph{grades} (finishing and conformation) are estimates, not
settled predictions. This follows from the project's oracle experiment, which
tested whether conformation is recoverable from a single 2D image:

\begin{caution}[title=The oracle result]
Shape descriptors were computed on the \emph{perfect manual annotation} (the best
mask any model could produce). \textbf{Zero of 384 statistical tests survived
multiple-comparison correction} at n=22. Even with perfect segmentation,
conformation is not reliably recoverable from a single 2D image at this sample
size. The bottleneck is the acquisition modality, not the model.
\end{caution}

The software therefore presents these grades as \pillalert{estimates}, not as
facts, and marks them in amber wherever they appear.

\section{Path to validation}

The validation the grades lack requires grades correctly paired to images, on a
larger sample. This is the data that Apex collects. As the paired dataset grows,
the \textbf{Physical correlation} tab (Station~7) allows testing whether the
surface measurements predict the physical reference and the assessor grades. This
addresses a step the earlier research could not complete.

\section{Modality caveats}

\begin{itemize}[leftmargin=1.3em,itemsep=3pt]
  \item \textbf{Fat cover (finishing)} is a depth property; a single surface
        image captures it only partly. Depth capture (Kinect) is supported in the
        schema for future work.
  \item \textbf{Standardised capture} affects every downstream measurement: a
        fixed background, uniform lighting and a consistent view improve the
        results. The software imposes a required view and offers background
        capture, but the physical setup is the operator's responsibility.
  \item \textbf{Sample size.} The current models come from n=22. Their estimates
        should be read as exploratory until the target sample is collected and
        the validation loop is run.
\end{itemize}

\begin{note}[title=Summary]
The earlier study was limited by data integrity rather than by its methods. The
software addresses this by collecting a reliable, paired dataset, so that
validation can follow. The limitations above account for the design decisions
described in this manual.
\end{note}
